\section{الفصل~\ref{sec:acet}. نضيف \nt{ACET}}

أفعال الأسيتيل كولين الثلاثة مقيسة في الشرائح وفي الدماغ الحي، والتعارض بين الكتابة والقراءة مبيَّن على نموذج الكتاب، أما أن \nt{ACET} يؤدي دور خط تمكين الكتابة ويخبر بعدم اليقين المتوقع ففرضية لها سند جزئي من التجربة.

{\footnotesize
\setlength{\tabcolsep}{3pt}\renewcommand{\arraystretch}{1.15}
\begin{longtable}{>{\raggedright}p{0.5cm}>{\raggedright}p{4.0cm}>{\raggedright}p{5.6cm}>{\raggedright}p{2.3cm}>{\raggedright\arraybackslash}p{1.7cm}}
\toprule
م & القضية & التجربة وما قيس & المصدر & الحالة \\
\midrule\endfirsthead
\toprule
م & القضية & التجربة وما قيس & المصدر & الحالة \\
\midrule\endhead
1 & عصبونات \nt{ACET} في الدماغ قليلة، مجموعة في بضع مجموعات، ومخارجها تتفرق في القشرة كلها: قناة بثّ & صبغ بأضداد لإنزيم تخليق الأسيتيل كولين ووسم راجع عند الجرذ: القشرة والحُصين واللوزة والبصلة الشمية والمهاد تتلقى الأسيتيل كولين أساسًا من ست مجموعات خلوية متراصة في الدماغ الأمامي القاعدي وفي الجذع & \cite{Mesulam1983} & مُقاس \\
2 & مستوى الأسيتيل كولين في القشرة مرتفع في اليقظة ومنخفض في النوم البطيء & غسيل مجهري في القشرة والحُصين عند قطط تتحرك بحرية مع تسجيل \foreignlanguage{english}{EEG}: تحرير الأسيتيل كولين في اليقظة الهادئة أعلى بـ$\sim100\%$، وفي اليقظة النشطة بـ$\sim175\%$، منه في النوم البطيء؛ وفي النوم السريع يكون في القشرة كما في اليقظة & \cite{Marrosu1995,Hasselmo1999} & مُقاس \\
3 & معنى الإشارة يحدده المستقبل: للأسيتيل كولين مستقبلات-قنوات سريعة، ومستقبلات بطيئة تطلق تفاعلات في الخلية (القضية~\ref{prop:receptor}) & مراجعة للقياسات: أثر الأسيتيل كولين في الاستثارية والنقل واللدونة يعتمد على موضع التحرير، ونوع المستقبل الفرعي (النيكوتينية قنوات أيونية، والمسكارينية عبر بروتينات \foreignlanguage{english}{G})، والخلية المتلقية & \cite{Picciotto2012} & مُقاس \\
4 & الفعل الأول: إضعاف الوصلات الداخلية دون المساس تقريبًا بالمداخل الخارجية (العامل $1-a$ في~\eqref{eq:acetnet}) & شرائح القشرة الشمية عند الجرذ: عند $100\,\mu\text{M}$ من ناهضات الأسيتيل كولين تنخفض كمونات الوصلات الداخلية (الطبقة \foreignlanguage{english}{Ib}) بأكثر من $60\%$، وكمونات الدخل الخارجي (الطبقة \foreignlanguage{english}{Ia}) بأقل من $15\%$، في الخلية نفسها؛ والكبت قبل مشبكي. شرائح الحُصين (\foreignlanguage{english}{CA1}): $89\%$ مقابل $40\%$ للمسار الداخلي ومسار الدخل & \cite{HasselmoBower1992,HasselmoSchnell1994} & مُقاس \\
5 & الفعل الثاني: تقوية الدخل (العامل $\frac{1+a}{2}$) & شرائح القشرة الحديثة: المستقبلات النيكوتينية تقوّي مشابك الدخل من المهاد ولا تؤثر في المشابك داخل القشرة؛ والمسكارينية تضعف النوعين. والأسيتيل كولين يرفع استثارية العصبونات (مراجعة) & \cite{Gil1997,Hasselmo2006} & مُقاس \\
6 & الفعل الثالث: تيسير تغيّر الأوزان (المعدل $\eta a$) & جرذ: التحفيز الكهربائي للنواة القاعدية (مصدر الأسيتيل كولين للقشرة) مقرونًا بنغمة يُحدث عند الحيوان البالغ إعادة تنظيم قوية لخريطة القشرة السمعية نحو الصوت المعروض & \cite{KilgardMerzenich1998,Hasselmo2006} & مُقاس \\
7 & قاعدة هيب للأنماط المتناثرة في المعادلة الثانية من~\eqref{eq:acetnet} تعطي ذاكرة ترابطية تُكمل النمط من جزئه & حساب سعة شبكة بأنماط متناثرة وقاعدة تغايرية: عند نسبة صغيرة $p$ من العصبونات النشطة تحفظ الشبكة أنماطًا أكثر بكثير مما عند $p=1/2$ & \cite{Tsodyks1988} & نظرية \\
8 & الكتابة عند $a$ منخفض تفسد الذاكرة: تكتب الشبكة ما تذكّرته لا ما رأته؛ وعند $a=0.1$ ليس الفساد أضعف منه عند $a=0$ & \texttt{models/\allowbreak acet\_memory.py}: $200$ عصبون، وخمسة أنماط في كل منها $20$، ونمط جديد يتقاسم $10$ عصبونات مع أحدها، و$10$ محاكاة. عند $a=0$ لا يُحفظ النمط الجديد، ويجرّ القديم $5.9$ عصبونًا غريبًا من $10$؛ وعند $a=0.1$ يُتذكَّر النمط الجديد ($0.97$)، لكن الاثنين يندمجان: الجديد يجرّ $7.3$ عصبونًا غريبًا، والقديم $7.9$؛ وبدءًا من $a=0.2$ أقل من واحد & اشتقاق في الفصل & نموذج \\
9 & القضية~\ref{prop:writeread}: لا يوجد مستوى $a$ تنجح عنده الكتابة والقراءة بالقدر نفسه (الشكل~\ref{fig:acetsim}) & البرنامج نفسه، ومعدل التعلّم $\eta a$: يُحفظ النمط الجديد كاملًا في عرض واحد عند $a\ge0.9$، و$0.8$ منه عند $a=0.75$، ولا يُحفظ عند $a\le0.5$. القراءة من نصف النمط: $1.0$ عند $a\le0.2$، و$0.96$ عند $a=0.3$، و$0.04$ عند $a=0.4$ & اشتقاق في الفصل & نموذج \\
10 & في الدماغ يبدّل سلك بثّ واحد، هو \nt{ACET}، بين الكتابة والقراءة & الفكرة مأخوذة من نماذج بُنيت على قياسات في الشرائح. وأكثر التجارب مباشرة عند الإنسان: حاجب المستقبلات المسكارينية سكوبولامين يسيء حفظ أزواج كلمات جديدة، وأشد ما يكون ذلك مع الأزواج المتداخلة مع القديمة، لكنه لا يسيء تذكّر الأزواج المحفوظة. لا يوجد قياس مباشر لوضعي الشبكة & \cite{HasselmoAndersonBower1992,Atri2004} & فرضية \\
11 & المرشح~\eqref{eq:kalman} أمثل؛ ووزن الدخل ومعدل التعلّم مقدار واحد $P/R$، وفي الحالة المستقرة $P=\sqrt{qR}$ والذاكرة $\sqrt{R/q}$ (القضية~\ref{prop:kalman}) & لا يتطلب تجربة: أمثلية مرشح كالمان-بيوسي نتيجة كلاسيكية في نظرية التقدير؛ والحالة المستقرة مشتقة في الفصل & \cite{KalmanBucy1961} & نظرية \\
12 & \nt{ACET} يخبر الشبكة بمقدار شبيه بـ$P$، أي عدم اليقين المتوقع & اقتُرح في نموذج احتمالي للانتباه. لا يوجد قياس مباشر لكون مستوى الأسيتيل كولين يتبع عدم اليقين في التقدير & \cite{YuDayan2005} & فرضية \\
13 & جزء من عصبونات \nt{ACET} يستجيب للمكافأة والعقاب في نحو عشرين ميلي ثانية، وبقوة تزداد كلما كان التعزيز أقل توقعًا & فأر، عصبونات كولينية في الدماغ الأمامي القاعدي معرَّفة بصريًا جينيًا: استجابة للمكافأة والعقاب بعد $18\pm3\ms$، يزداد حجمها مع المفاجأة في التعزيز، والاستجابات متشابهة في عصبونات نواتين & \cite{Hangya2015} & مُقاس \\
14 & \nt{NORA} يلاحظ التغيّر غير المتوقع في العالم، و\nt{ACET} يُبقي الشبكة في وضع الكتابة & تقسيم العمل مقترح في نموذج. وبصورة غير مباشرة: عند الإنسان يتتبع قطر الحدقة المرتبط بـ\nt{NORA} التغيرات ومعدل التعلّم. لا يوجد اختبار مباشر للزوج \foreignlanguage{english}{\nt{NORA}--\nt{ACET}} & \cite{YuDayan2005,Nassar2012} & فرضية \\
15 & في النوم البطيء تعيد الشبكة في وضع القراءة تشغيل ما كُتب نهارًا، وهذه الإعادات تنقله إلى الذاكرة الطويلة الأمد & تسجيلات مجموعات عصبونات الحُصين عند الجرذ: الخلايا التي عملت معًا نهارًا تنطلق معًا أكثر في النوم البطيء التالي، وبالترتيب نفسه، وهذا مقيس. وللنقل: كبت تموجات الحُصين بعد التعلّم يسيء الذاكرة المكانية، وإطالة التموجات التلقائية يحسّنها؛ أما أن السبب هو المستوى المنخفض لـ\nt{ACET} فغير مبرهن & \cite{WilsonMcNaughton1994,Skaggs1996,Girardeau2009,FernandezRuiz2019,Hasselmo1999} & فرضية \\
\bottomrule
\end{longtable}}
